\documentclass[../../../../main.tex]{subfiles}
\begin{document}

\begin{flushright}
\footnotesize\textit{【自動文字起こし・要確認】}
\end{flushright}

[解] $0 \le a < \pi/4, \ 0 \le b < \pi/4 \cdots$ ①

[補題1] $0 \le x < \pi/4$ において、$f(x) = \tan x$ は下に凸 \\
$f'(x) = \frac{1}{\cos^2 x} > 0$, $f''(x) = 2 \frac{\tan x}{\cos^2 x} > 0$ から示された。$\square$

[補題2] $0 < x < \pi/4$ において $g(x) = \log(\tan x)$ は上に凸 \\
$g'(x) = \frac{1}{\cos^2 x \cdot \tan x} > 0$, $g''(x) = - \frac{1}{\sin^2 x} + \frac{1}{\cos^2 x} = \frac{2\sin^2 x - 1}{\sin^2 x \cos^2 x} < 0$ \\
から示された。$\square$

・[補題1]から凸不等式より右側の不等式は示された。$\square$

次に左側について \\
$1^\circ \ a>0, b>0$ の時 \\
両辺 $\log$ をとることができて[補題2]から凸不等式から示された。

$2^\circ \ ab=0$ の時
\[
\begin{cases}
\sqrt{\tan 0 \cdot \tan b} = 0 \\
\tan\left(\frac{0+b}{2}\right) \ge 0
\end{cases}
\]
から明らかに成立

より、左側も示された。$\square$

\end{document}
